\section*{Наглядное введение}

Метод геометрии масс основан на простой и наглядной идее: геометрические точки можно мысленно рассматривать как \textit{материальные точки}, то есть точки, в которых сосредоточена некоторая масса.

Если есть две точки $A$ и $B$ с массами $m_1$ и $m_2$, то им можно сопоставить одну новую точку --- \textit{центр масс}. Эта точка лежит на отрезке $AB$ и расположена так, что более тяжёлая точка оказывается к ней ближе. Иначе говоря, центр масс делит отрезок в отношении, обратном массам:
\[
m_1 \cdot ZA = m_2 \cdot ZB.
\]

Это удобно понимать так: две точки с массами можно заменить одной точкой с общей массой
\[
m_1 + m_2,
\]
не меняя <<равновесия>> системы.

Та же идея работает и для нескольких точек. Если дано много точек с массами, то их можно последовательно объединять: сначала найти центр масс двух точек, потом добавить к нему третью точку, и так далее. В результате вся система заменяется одной точкой --- её центром масс.

Главная ценность метода в том, что геометрическая задача после такой <<загрузки массами>> часто становится гораздо понятнее. Вместо сложных рассуждений можно использовать простые свойства центра масс: где он находится, как делит отрезки и как меняется при объединении точек.

Именно поэтому геометрия масс полезна при решении задач на отрезки, медианы, отношения и точки пересечения.

\section*{Математическое введение}

\textbf{Материальная точка.}
Материальной точкой $mA$ будем называть точку $A$, которой сопоставлена масса $m$.

\medskip

\textbf{Центр масс системы точек.}
Центром масс системы
\[
m_1A_1,\; m_2A_2,\; \dots,\; m_nA_n
\]
называется точка $Z$, для которой выполняется равенство
\[
m_1\overrightarrow{ZA_1}+m_2\overrightarrow{ZA_2}+\dots+m_n\overrightarrow{ZA_n}=\vec 0.
\]

\medskip

\textbf{Центр масс двух точек.}
Если система состоит из двух точек $m_1A$ и $m_2B$, то её центр масс $Z$ лежит на отрезке $AB$ и делит его в отношении, обратном массам:
\[
m_1\cdot ZA = m_2\cdot ZB.
\]
Значит,
\[
\frac{ZA}{ZB}=\frac{m_2}{m_1}.
\]

\medskip

\textbf{Формула центра масс относительно точки $O$.}
Если $Z$ --- центр масс системы
\[
m_1A_1,\; m_2A_2,\; \dots,\; m_nA_n,
\]
то для любой точки $O$ верно:
\[
\overrightarrow{OZ}=
\frac{m_1\overrightarrow{OA_1}+m_2\overrightarrow{OA_2}+\dots+m_n\overrightarrow{OA_n}}
{m_1+m_2+\dots+m_n}.
\]

\medskip

\textbf{Свойство объединения масс.}
Если несколько точек заменить их центром масс, а в этот центр поместить сумму этих масс, то центр масс всей системы не изменится.

\subsection*{Примеры решения задач}

\textbf{Пример 1.}
На отрезке $AB$ дана точка $M$, причём
\[
AM:MB=2:3.
\]
Какие массы надо поместить в точки $A$ и $B$, чтобы точка $M$ была центром масс?

\medskip

\textbf{Решение.}
Пусть в точках $A$ и $B$ стоят массы $m_A$ и $m_B$.
Так как $M$ --- центр масс двух точек, имеем:
\[
m_A\cdot AM=m_B\cdot BM.
\]
Подставим отношение:
\[
m_A\cdot 2 = m_B\cdot 3.
\]
Отсюда
\[
m_A:m_B=3:2.
\]


\bigskip

\textbf{Пример 2.}
Докажем, что три медианы треугольника пересекаются в одной точке, причём эта точка делит каждую медиану в отношении $2:1$, считая от вершины.

\medskip

\textbf{Решение.}
Рассмотрим треугольник $ABC$ и поместим в его вершины равные массы:
\[
m_A=m_B=m_C=1.
\]
Пусть $A_1$ --- середина стороны $BC$.
Тогда точки $B$ и $C$ можно заменить их центром масс, то есть точкой $A_1$, в которой будет сосредоточена масса
\[
1+1=2.
\]
Значит, центр масс всего треугольника лежит на прямой $AA_1$, то есть на медиане.

Точно так же, если $B_1$ и $C_1$ --- середины сторон $AC$ и $AB$, то центр масс лежит и на медианах $BB_1$ и $CC_1$.

Следовательно, все три медианы проходят через одну точку $Z$ --- центр масс системы.

Теперь найдём, как эта точка делит медиану $AA_1$.
В точке $A$ стоит масса $1$, а в точке $A_1$ --- масса $2$.
По правилу для двух точек:
\[
1\cdot ZA = 2\cdot ZA_1.
\]
Значит,
\[
ZA:ZA_1=2:1.
\]
Следовательно,
\[
AZ:ZA_1=2:1.
\]

\bigskip

\textbf{Пример 3.}
В вершинах параллелограмма $ABCD$ помещены массы
\[
p,\; q,\; p,\; q
\]
соответственно. Докажем, что центр масс этой системы находится в точке пересечения диагоналей.

\medskip

\textbf{Решение.}
Так как в вершинах $A$ и $C$ стоят равные массы $p$ и $p$, их центр масс --- середина диагонали $AC$.
Обозначим её через $O$.
В точке $O$ тогда можно считать сосредоточенной массу
\[
2p.
\]

Аналогично, точки $B$ и $D$ с равными массами $q$ и $q$ можно заменить их центром масс --- серединой диагонали $BD$. 

Но в параллелограмме диагонали пересекаются и делятся пополам, значит середина $AC$ совпадает с серединой $BD$.
Это и есть точка $O$. Следовательно, после объединения масс вся система сводится к одной точке $O$, а значит $O$ и есть центр масс.